\documentclass[11pt]{article}
\newcommand{\ListaTitulo}{Gabarito --- Lista 07}
% Preâmbulo compartilhado — MATD48, Listas de Exercícios 2026
% Usado por Lista{NN}.tex e Gabarito{NN}.tex via % Preâmbulo compartilhado — MATD48, Listas de Exercícios 2026
% Usado por Lista{NN}.tex e Gabarito{NN}.tex via % Preâmbulo compartilhado — MATD48, Listas de Exercícios 2026
% Usado por Lista{NN}.tex e Gabarito{NN}.tex via \input{preamble}

\usepackage[utf8]{inputenc}
\usepackage[T1]{fontenc}
\usepackage[brazil]{babel}
\usepackage{fullpage}
\usepackage{amsmath,amssymb}
\usepackage{enumitem}
\usepackage{xcolor}
\usepackage{fancyhdr}
\usepackage{titlesec}
\usepackage{listings}
\usepackage{hyperref}
\usepackage{booktabs}
\usepackage{graphicx}

\definecolor{matd48azul}{HTML}{0B4F6C}
\definecolor{matd48verde}{HTML}{2E8B57}

\titleformat{\section}{\normalfont\Large\bfseries\color{matd48azul}}{\thesection}{1em}{}
\titleformat{\subsection}{\normalfont\large\bfseries\color{matd48azul}}{\thesubsection}{1em}{}

\providecommand{\ListaTitulo}{Lista de Exercícios}

\setlength{\headheight}{14pt}
\pagestyle{fancy}
\fancyhf{}
\lhead{\small MATD48 --- Planejamento de Experimentos A}
\rhead{\small \ListaTitulo}
\cfoot{\thepage}
\renewcommand{\headrulewidth}{0.4pt}

\lstset{
  basicstyle=\ttfamily\small,
  breaklines=true,
  frame=single,
  columns=fullflexible,
  backgroundcolor=\color{gray!8},
  commentstyle=\color{matd48verde},
  keywordstyle=\color{matd48azul}\bfseries,
  language=R,
  extendedchars=true,
  literate=%
    {á}{{\'a}}1 {à}{{\`a}}1 {â}{{\^a}}1 {ã}{{\~a}}1
    {é}{{\'e}}1 {ê}{{\^e}}1 {í}{{\'i}}1 {ó}{{\'o}}1
    {ô}{{\^o}}1 {õ}{{\~o}}1 {ú}{{\'u}}1 {ç}{{\c c}}1
    {Á}{{\'A}}1 {À}{{\`A}}1 {Â}{{\^A}}1 {Ã}{{\~A}}1
    {É}{{\'E}}1 {Ê}{{\^E}}1 {Í}{{\'I}}1 {Ó}{{\'O}}1
    {Ô}{{\^O}}1 {Õ}{{\~O}}1 {Ú}{{\'U}}1 {Ç}{{\c C}}1
}

\hypersetup{colorlinks=true, linkcolor=matd48azul, urlcolor=matd48azul, citecolor=matd48azul}

\newcommand{\cabecalho}[2]{%
  \begin{center}
    {\Large\bfseries\color{matd48azul} #1}\\[0.3em]
    {\large #2}\\[0.3em]
    \rule{\linewidth}{0.6pt}
  \end{center}
  \vspace{0.5em}
}

\newenvironment{questao}[1]{\vspace{1em}\noindent\textbf{Questão #1.}\hspace{0.4em}}{\par}
\newenvironment{solucao}{\vspace{0.3em}\noindent\textit{Solução.}\hspace{0.4em}\color{black!85}}{\par\vspace{0.5em}}


\usepackage[utf8]{inputenc}
\usepackage[T1]{fontenc}
\usepackage[brazil]{babel}
\usepackage{fullpage}
\usepackage{amsmath,amssymb}
\usepackage{enumitem}
\usepackage{xcolor}
\usepackage{fancyhdr}
\usepackage{titlesec}
\usepackage{listings}
\usepackage{hyperref}
\usepackage{booktabs}
\usepackage{graphicx}

\definecolor{matd48azul}{HTML}{0B4F6C}
\definecolor{matd48verde}{HTML}{2E8B57}

\titleformat{\section}{\normalfont\Large\bfseries\color{matd48azul}}{\thesection}{1em}{}
\titleformat{\subsection}{\normalfont\large\bfseries\color{matd48azul}}{\thesubsection}{1em}{}

\providecommand{\ListaTitulo}{Lista de Exercícios}

\setlength{\headheight}{14pt}
\pagestyle{fancy}
\fancyhf{}
\lhead{\small MATD48 --- Planejamento de Experimentos A}
\rhead{\small \ListaTitulo}
\cfoot{\thepage}
\renewcommand{\headrulewidth}{0.4pt}

\lstset{
  basicstyle=\ttfamily\small,
  breaklines=true,
  frame=single,
  columns=fullflexible,
  backgroundcolor=\color{gray!8},
  commentstyle=\color{matd48verde},
  keywordstyle=\color{matd48azul}\bfseries,
  language=R,
  extendedchars=true,
  literate=%
    {á}{{\'a}}1 {à}{{\`a}}1 {â}{{\^a}}1 {ã}{{\~a}}1
    {é}{{\'e}}1 {ê}{{\^e}}1 {í}{{\'i}}1 {ó}{{\'o}}1
    {ô}{{\^o}}1 {õ}{{\~o}}1 {ú}{{\'u}}1 {ç}{{\c c}}1
    {Á}{{\'A}}1 {À}{{\`A}}1 {Â}{{\^A}}1 {Ã}{{\~A}}1
    {É}{{\'E}}1 {Ê}{{\^E}}1 {Í}{{\'I}}1 {Ó}{{\'O}}1
    {Ô}{{\^O}}1 {Õ}{{\~O}}1 {Ú}{{\'U}}1 {Ç}{{\c C}}1
}

\hypersetup{colorlinks=true, linkcolor=matd48azul, urlcolor=matd48azul, citecolor=matd48azul}

\newcommand{\cabecalho}[2]{%
  \begin{center}
    {\Large\bfseries\color{matd48azul} #1}\\[0.3em]
    {\large #2}\\[0.3em]
    \rule{\linewidth}{0.6pt}
  \end{center}
  \vspace{0.5em}
}

\newenvironment{questao}[1]{\vspace{1em}\noindent\textbf{Questão #1.}\hspace{0.4em}}{\par}
\newenvironment{solucao}{\vspace{0.3em}\noindent\textit{Solução.}\hspace{0.4em}\color{black!85}}{\par\vspace{0.5em}}


\usepackage[utf8]{inputenc}
\usepackage[T1]{fontenc}
\usepackage[brazil]{babel}
\usepackage{fullpage}
\usepackage{amsmath,amssymb}
\usepackage{enumitem}
\usepackage{xcolor}
\usepackage{fancyhdr}
\usepackage{titlesec}
\usepackage{listings}
\usepackage{hyperref}
\usepackage{booktabs}
\usepackage{graphicx}

\definecolor{matd48azul}{HTML}{0B4F6C}
\definecolor{matd48verde}{HTML}{2E8B57}

\titleformat{\section}{\normalfont\Large\bfseries\color{matd48azul}}{\thesection}{1em}{}
\titleformat{\subsection}{\normalfont\large\bfseries\color{matd48azul}}{\thesubsection}{1em}{}

\providecommand{\ListaTitulo}{Lista de Exercícios}

\setlength{\headheight}{14pt}
\pagestyle{fancy}
\fancyhf{}
\lhead{\small MATD48 --- Planejamento de Experimentos A}
\rhead{\small \ListaTitulo}
\cfoot{\thepage}
\renewcommand{\headrulewidth}{0.4pt}

\lstset{
  basicstyle=\ttfamily\small,
  breaklines=true,
  frame=single,
  columns=fullflexible,
  backgroundcolor=\color{gray!8},
  commentstyle=\color{matd48verde},
  keywordstyle=\color{matd48azul}\bfseries,
  language=R,
  extendedchars=true,
  literate=%
    {á}{{\'a}}1 {à}{{\`a}}1 {â}{{\^a}}1 {ã}{{\~a}}1
    {é}{{\'e}}1 {ê}{{\^e}}1 {í}{{\'i}}1 {ó}{{\'o}}1
    {ô}{{\^o}}1 {õ}{{\~o}}1 {ú}{{\'u}}1 {ç}{{\c c}}1
    {Á}{{\'A}}1 {À}{{\`A}}1 {Â}{{\^A}}1 {Ã}{{\~A}}1
    {É}{{\'E}}1 {Ê}{{\^E}}1 {Í}{{\'I}}1 {Ó}{{\'O}}1
    {Ô}{{\^O}}1 {Õ}{{\~O}}1 {Ú}{{\'U}}1 {Ç}{{\c C}}1
}

\hypersetup{colorlinks=true, linkcolor=matd48azul, urlcolor=matd48azul, citecolor=matd48azul}

\newcommand{\cabecalho}[2]{%
  \begin{center}
    {\Large\bfseries\color{matd48azul} #1}\\[0.3em]
    {\large #2}\\[0.3em]
    \rule{\linewidth}{0.6pt}
  \end{center}
  \vspace{0.5em}
}

\newenvironment{questao}[1]{\vspace{1em}\noindent\textbf{Questão #1.}\hspace{0.4em}}{\par}
\newenvironment{solucao}{\vspace{0.3em}\noindent\textit{Solução.}\hspace{0.4em}\color{black!85}}{\par\vspace{0.5em}}


\begin{document}

\cabecalho{Gabarito --- Lista de Exercícios 07}{Contrastes, comparações múltiplas e polinômios ortogonais (Capítulo 3, segunda metade / Aula 07)}

\begin{questao}{1} \textbf{(Ciência de dados --- cálculo de um contraste)}
\end{questao}
\begin{solucao}
\begin{enumerate}[label=(\alph*)]
  \item $\boldsymbol\lambda = (3, -1, -1, -1)'$. Soma: $3 - 1 - 1 - 1 = 0$, logo
    $\mathbf 1'\boldsymbol\lambda = 0$: é de fato um contraste.
  \item $\hat L = 3(860) - 940 - 1080 - 1120 = 2580 - 3140 = -560$.
  \item $\sum_i c_i^2 = 9+1+1+1 = 12$, $n=10$, logo
    \[
    SQ_L = \frac{\hat L^2}{\sum_i c_i^2/n} = \frac{(-560)^2}{12/10} = \frac{313\,600}{1{,}2} \approx
    261\,333{,}3.
    \]
    $F = SQ_L/QME = 261\,333{,}3/4200 \approx 62{,}22$. Como $62{,}22 > 4{,}11 =
    F_{0{,}05;\,1,36}$, \textbf{rejeitamos $H_0: L=0$}: usar qualquer script de terceiros
    aumenta o tempo de carregamento em relação a não usar nenhum, com forte evidência
    estatística.
\end{enumerate}
\end{solucao}

\begin{questao}{2} \textbf{(Ciência de dados --- contrastes ortogonais e partição de $SQ_{\text{Trat}}$)}
\end{questao}
\begin{solucao}
\begin{enumerate}[label=(\alph*)]
  \item $\boldsymbol\lambda_1'\boldsymbol\lambda_2 = 3(0)+(-1)(2)+(-1)(-1)+(-1)(-1) = 0-2+1+1=0$.
    $\boldsymbol\lambda_1'\boldsymbol\lambda_3 = 3(0)+(-1)(0)+(-1)(1)+(-1)(-1) = 0+0-1+1=0$.
    $\boldsymbol\lambda_2'\boldsymbol\lambda_3 = 0(0)+2(0)+(-1)(1)+(-1)(-1) = 0+0-1+1=0$. Os três
    produtos internos são nulos: os três contrastes são mutuamente ortogonais.
  \item $\hat L_2 = 2(940) - 1080 - 1120 = 1880-2200=-320$. $\sum c_i^2 = 4+1+1=6$, logo
    $SQ_{L_2} = 320^2/(6/10) = 102\,400/0{,}6 \approx 170\,666{,}7$.
    $\hat L_3 = 1080-1120=-40$. $\sum c_i^2=1+1=2$, logo $SQ_{L_3}=40^2/(2/10)=1600/0{,}2=8000$.
  \item $SQ_{L_1}+SQ_{L_2}+SQ_{L_3} = 261\,333{,}3 + 170\,666{,}7 + 8000 = 440\,000 = SQ_{\text{Trat}}$.
    A partição se confirma exatamente.
  \item $F_1 \approx 62{,}2$ e $F_2 = 170\,666{,}7/4200 \approx 40{,}6$ são ambos muito maiores
    que $F_{0{,}05;1,36}\approx 4{,}11$: significativos. $F_3 = 8000/4200 \approx 1{,}90 < 4{,}11$:
    não significativo. Conclusão: a diferença entre configurações vem de "nenhum script vs.
    demais" e de "Analytics vs. (Anúncios, Combo)"; Anúncios e Combo, especificamente, não
    diferem entre si de forma detectável neste experimento.
\end{enumerate}
\end{solucao}

\begin{questao}{3} \textbf{(Psicologia --- Scheffé vs. teste ingênuo)}
\end{questao}
\begin{solucao}
\begin{enumerate}[label=(\alph*)]
  \item $F = 185/40 = 4{,}625$. Como $4{,}625 > 4{,}12 = F_{0{,}05;1,35}$, o teste ingênuo (não
    corrigido) rejeitaria $H_0$: o contraste "pareceria" significativo.
  \item Valor crítico de Scheffé: $(t-1)F_{0{,}05;4,35} = 4 \times 2{,}64 = 10{,}56$. Como
    $F = 4{,}625 < 10{,}56$, \textbf{não rejeitamos $H_0$ pelo critério de Scheffé}: o contraste
    deixa de ser significativo quando corrigido para o fato de ter sido escolhido a posteriori.
  \item O teste do item (a) usa $\alpha = 0{,}05$ como se este fosse o \emph{único} contraste
    possível a ser testado, definido antes de qualquer olhar aos dados. Mas o pesquisador, ao
    examinar os dados antes de escolher qual comparação testar, efetivamente considerou (mesmo
    que implicitamente) um número muito maior de contrastes possíveis, e escolheu justamente
    aquele que parecia mais promissor. Isso é análogo ao problema de comparações múltiplas
    (ver a seção sobre procedimentos de comparação múltipla no livro-texto, Capítulo 3): a
    probabilidade de encontrar \emph{algum}
    contraste "significativo" por puro acaso, entre a enorme quantidade de contrastes possíveis
    com $t=5$ grupos, é muito maior que 5\%. O critério de Scheffé é desenhado exatamente para
    controlar essa taxa de erro por família sobre \emph{todos} os contrastes possíveis
    simultaneamente, não apenas sobre o único que acabou sendo reportado.
\end{enumerate}
\end{solucao}

\begin{questao}{4} \textbf{(Dedução --- independência de contrastes ortogonais)}
\end{questao}
\begin{solucao}
\begin{enumerate}[label=(\alph*)]
  \item Como as $t$ unidades experimentais em grupos diferentes são independentes, as médias
    amostrais $\bar y_{1\cdot},\ldots,\bar y_{t\cdot}$ são independentes entre si, com
    $\mathrm{Var}(\bar y_{i\cdot}) = \sigma^2/n$ para todo $i$ (balanceado). Logo
    $\mathrm{Cov}(\bar{\mathbf y}) = (\sigma^2/n)\mathbf I_t$. Para combinações lineares
    $\hat L_1 = \boldsymbol\lambda_1'\bar{\mathbf y}$ e $\hat L_2 = \boldsymbol\lambda_2'\bar{\mathbf y}$,
    \[
    \mathrm{Cov}(\hat L_1,\hat L_2) = \boldsymbol\lambda_1'\,\mathrm{Cov}(\bar{\mathbf y})\,
    \boldsymbol\lambda_2 = \boldsymbol\lambda_1'\left(\frac{\sigma^2}{n}\mathbf I_t\right)
    \boldsymbol\lambda_2 = \frac{\sigma^2}{n}\boldsymbol\lambda_1'\boldsymbol\lambda_2.
    \]
  \item Se $\boldsymbol\lambda_1'\boldsymbol\lambda_2=0$ (ortogonalidade, por definição), então
    imediatamente $\mathrm{Cov}(\hat L_1,\hat L_2) = (\sigma^2/n)\cdot 0 = 0$.
  \item $\bar{\mathbf y}$ é um vetor conjuntamente normal (combinação linear de $Y_{ij}$ normais
    independentes). $\hat L_1$ e $\hat L_2$ são ambas combinações lineares de $\bar{\mathbf y}$,
    logo o par $(\hat L_1, \hat L_2)$ também é conjuntamente normal (bivariado). Para vetores
    conjuntamente normais, é um resultado padrão que covariância nula implica independência
    estatística (isso \textbf{não} vale, em geral, para variáveis não normais — a normalidade é
    essencial aqui). Logo $\hat L_1 \perp \hat L_2$. Como $SQ_{L_1} = \hat L_1^2/(\sum c_i^2/n)$ é
    função determinística apenas de $\hat L_1$, e $SQ_{L_2}$ apenas de $\hat L_2$, a independência
    de $\hat L_1$ e $\hat L_2$ implica a independência de $SQ_{L_1}$ e $SQ_{L_2}$ (uma função
    mensurável de uma variável independente de outra função mensurável de outra variável
    preserva independência).
\end{enumerate}
\end{solucao}

\begin{questao}{5} \textbf{(Dedução --- coeficientes do polinômio ortogonal linear, $t=3$)}
\end{questao}
\begin{solucao}
\begin{enumerate}[label=(\alph*)]
  \item Os níveis codificados são $-1, 0, 1$. Proporcionalidade direta ($c_i = k\cdot
    \text{nível}_i$) com $k=1$ dá $\boldsymbol\lambda_L = (-1, 0, 1)$ --- já são inteiros de
    módulo mínimo (o próximo múltiplo inteiro, $k=2$, daria $(-2,0,2)$, com módulo maior). A
    condição (i) é automática: como os níveis codificados $-1,0,1$ já somam zero por construção
    (são simétricos em torno do centro), qualquer múltiplo escalar deles também soma zero:
    $\sum_i c_i = k\sum_i \text{nível}_i = k \cdot 0 = 0$.
  \item $\boldsymbol\lambda_L'\boldsymbol\lambda_Q = (-1)(1) + (0)(-2) + (1)(1) = -1+0+1=0$:
    são ortogonais.
  \item Para $t=3$ níveis, o número máximo de contrastes mutuamente ortogonais é $t-1=2$.
    $\boldsymbol\lambda_L$ e $\boldsymbol\lambda_Q$ são exatamente 2 contrastes mutuamente
    ortogonais (item b), logo esgotam a contagem máxima: juntos, eles decompõem completamente
    $SQ_{\text{Trat}}$ (2 graus de liberdade), sem espaço para nenhum grau polinomial adicional
    (não existe componente cúbico com apenas 3 níveis).
\end{enumerate}
\end{solucao}

\begin{questao}{6} \textbf{(Agricultura --- interpretação de polinômios ortogonais)}
\end{questao}
\begin{solucao}
\begin{enumerate}[label=(\alph*)]
  \item $705 + 98 + 9 = 812$, que é exatamente o $SQ$ de "Dose (total)" --- a partição se
    confirma.
  \item Ao nível de 5\%, os componentes \textbf{linear} ($p<0{,}001$) e \textbf{quadrático}
    ($p=0{,}028$) são significativos; o \textbf{cúbico} ($p=0{,}491$) não é. Isso indica que a
    produtividade cresce com a dose de nitrogênio de forma predominantemente linear, mas com uma
    curvatura estatisticamente detectável (componente quadrático) --- tipicamente, um padrão de
    retornos decrescentes: cada incremento de dose aumenta a produtividade, mas em magnitude
    progressivamente menor à medida que a dose cresce, sem evidência de uma reversão adicional
    (S ou oscilação) que um termo cúbico capturaria.
  \item Polinômios ortogonais exigem níveis \textbf{quantitativos, ordenados e igualmente
    espaçados} para que os coeficientes tabulados representem tendência (linear, quadrática...)
    de forma interpretável. Variedades de milho são um fator \textbf{categórico (nominal)}: não
    há uma noção de "meia unidade" entre duas variedades, nem uma ordem numérica natural entre
    elas, de modo que "tendência linear entre variedades" não tem significado --- para
    variedades, o procedimento adequado é comparação múltipla (Tukey, Dunnett etc.) ou contrastes
    específicos definidos pelo interesse do pesquisador, não uma decomposição polinomial.
\end{enumerate}
\end{solucao}

\end{document}
